\scanpage{93}
\mspage{92}
\mspage{42}

d'où un homomorphisme $f^{\unsure{*}}(G)\longrightarrow F'$ et ce dernier est évidemment compatible avec les données de descente. Il reste à examiner si cet homomorphisme est un \uline{isomorphisme.}

\figure{\begin{tikzpicture}\draw (0,0) .. controls (-0.2,0.15) and (-0.15,0.3) .. (0.1,0.3) -- (4.1,0.48) .. controls (4.35,0.48) and (4.35,0.65) .. (4.2,0.73);\end{tikzpicture}}

Noter que $G$ est aussi définissable par l'exactitude de
\diagram{\begin{tikzcd}G\arrow[r]&f_{\unsure{*}}(F')\arrow[r,shift left=.65ex]\arrow[r,shift right=.65ex]&g_{\unsure{*}}(F'')\end{tikzcd}},
donc pour tout changement de base $Y\xrightarrow{h}X$ \uline{qui commute à la formation de} $f_{\unsure{*}}(F')$ \uline{et} $g_{\unsure{*}}(F'')$, $h^{\unsure{*}}(G)=G_Y$ s'identifie au faisceau ``descendu'' de $F'_{Y'}$ par $Y'\longrightarrow Y$, et bien entendu l'homomorphisme $f_Y^{\unsure{*}}(G_Y)\longrightarrow(F_Y)'$ est celui déduit de $f^{\unsure{*}}(G)\longrightarrow F'$ par changement de base. Or supposons que \uline{le changement de base} $f:Y=X'\longrightarrow X$ \uline{commute à la formation de} $f_{\unsure{*}}(F')$ \uline{et} $g_{\unsure{*}}(F'')$, et notons que comme $Y'=X'\inserted{\times}_X X'$ a une section sur $Y'=X'$, donc $Y'\longrightarrow Y$ est un morphisme de descente \uline{effective} pour toute catégorie fibrée, en particulier pour $\unsure{\mathfrak F|\mathcal F}$, il s'ensuit que $f_Y^{\unsure{*}}(G_Y)\longrightarrow(F_Y)'$ est un isomorphisme, donc $f^{\unsure{*}}(G)\longrightarrow F'$ devient un isomorphisme après changement de base $Y'=X'\quad{}_X X'\longrightarrow X'$. Comme ce dernier à une section, $f^{\unsure{*}}(G)\longrightarrow F'$ est un isomorphisme, donc la donnée de descente envisagée sur $F'$ est effective.

9.4.2. Ceci prouve le théorème 9.4. dans le cas a), grâce à 5.6. Le cas b) résulte également du fait que si $f$ est propre, alors $f_{\unsure{*}}$ commute à \uline{tout} changement de base $Y'\longrightarrow Y$, (qui sera prouvé dans un exposé ultérieur).

9.4.3. Dans le cas c), la question étant locale sur $Y$ on peut
